\section{Consequences for ideal contraction budgets}
\label{sec:repair}

Let $y$ be the centered star curvature and consider the ideal contraction
$y\mapsto ay$, $0<a<1$.  The required displacement is
$\delta y=-(1-a)y$.  Theorem~\ref{thm:spill-spectrum} gives
\[
 \min_{R_\ell x=\delta y}\norm{B_\ell x}^2
 \ge \frac{37}{33}(1-a)^2\norm y^2.
\]
If the external curvature initially vanishes, the most favorable quadratic
action accounting therefore satisfies
\begin{equation}
\label{eq:action-budget}
 \Delta\mathcal A
 \ge \frac1{2t}\left[
   -(1-a^2)+\frac{37}{33}(1-a)^2
 \right]\norm y^2.
\end{equation}
The right-hand side may be negative, so action reduction remains possible.
However, its magnitude is strictly smaller than the star-only prediction
$-(1-a^2)\norm y^2/(2t)$.

In a change-of-variables argument, the $15$ real centered color modes of a
six-staple star also produce the ideal coordinate Jacobian $a^{15}$.  If $J$
collects every further nonnegative logarithmic loss and $\Gamma$ is the desired
free-energy gain, a sufficient algebraic design condition for a spill
coefficient $\lambda$ is
\begin{equation}
\label{eq:design-threshold}
 \lambda\le
 \frac{(1-a^2)D_0/(2t)-15\log(1/a)-J-\Gamma}
 {(1-a)^2D_0/(2t)}.
\end{equation}
Here $D_0$ is the minimum dispersion on the repaired event.  Substitution of
$\lambda=391/239$ tests the worst block mode, while $37/33$ tests the best.

Equation~\eqref{eq:design-threshold} is a design inequality, not a proved
nonlinear measure transformation.  In particular, the Jacobian of a genuine
coordinated link map need not equal the ideal staple-coordinate Jacobian.
This distinction is exactly the point: an operation that is transparent in
one representation acquires compatibility, spillover, and Jacobian costs when
realized in another.
